\section{Balance armónico y espectros locales}
\label{app:wolfram-master}

\subsection{Funciones descriptivas y obstrucciones de balance}

Con la convención \(W_+\), se obtiene
\begin{align}
 r^{\mathsf T}(zI-P)^{-1}b
 &=-\frac{\alpha Q(z)}{\Delta_\ell(z)},&
 \det(zI-P)&=\Delta_\ell(z),\\
 N_{\mathrm{cub}}(A)&=\frac34A^2,&
 N_{\mathrm{sat}}(A)
 &=\frac{2\Delta}{\pi}
 \left[
 \arcsin\frac1A+\frac{\sqrt{A^2-1}}{A^2}
 \right],\\
 N_{\arctan}(A)&=
 \frac{2a_2}{\rho A^2}
 \left(\sqrt{1+\rho^2A^2}-1\right).
\end{align}
La expresión cúbica corresponde a \(\psi(u)=u^3\), con \(A>0\); la de
saturación usa \(\Delta=m_0-m_1\) y se aplica a \(A>1\). Para la arctangente se supone
\(A>0\) y \(\rho>0\). Las dos últimas expresiones se obtienen de la
proyección \cref{eq:df_centrada}, mediante la integración por tramos y por
partes desarrollada junto a \cref{eq:df_saturacion,eq:df_arctan}.
En el caso cúbico,
\[
 N_{\mathrm{cub}}(A)=\frac{2A^2}{\pi}\int_0^\pi\sin^4\theta\,\mathrm d\theta
 =\frac{2A^2}{\pi}\frac{3\pi}{8}=\frac34A^2,
\]
pues \(\sin^4\theta=(3-4\cos2\theta+\cos4\theta)/8\).

Para Kalman--Fitts,
\[
 W(s)=-\frac{s^2}{p(s)},\qquad
 \omega_0^2=\frac{2\beta^2+m_1^2+m_2^2}{2},
\]
y
\[
 k=-\left[
 4\beta^2+
 \frac{(m_1^2-m_2^2)^2}
 {2(2\beta^2+m_1^2+m_2^2)}
 \right]
 =-0.0431687011573845.
\]
Para PLL,
\[
 \Im W_+(\mathrm i\omega)=
 \frac{L[1+\omega^2\tau_2(\tau_1+\tau_2)]}
 {2\omega[1+\omega^2(\tau_1+\tau_2)^2]}>0.
\]
Se supone \(L,\tau_1,\tau_2,\omega>0\); en la convención alternativa,
\(\Im W_-=-\Im W_+<0\). En ambas convenciones la parte imaginaria conserva
su signo y no se anula a frecuencia positiva.
La ganancia negativa de Kalman--Fitts y el signo de la parte imaginaria
del PLL obstruyen el cierre directo de función descriptiva en los casos
considerados.

\subsection{Espectro del Lorenz generalizado}

En el origen, el jacobiano se descompone en un bloque de dimensión dos y la
dirección \(z\). Por ello,
\[
 p_0(\lambda)=(\lambda+1)
 \bigl[(\lambda+3.4)(\lambda+1)-3.4\cdot6.8\bigr]
 =(\lambda+1)(\lambda^2+4.4\lambda-19.72).
\]
Las raíces del segundo factor son \(-2.2\pm\sqrt{24.56}\).
El espectro en el origen es
\[
 \{-7.15580467734555,\ 2.75580467734555,\ -1\},
\]
y en los dos equilibrios no nulos,
\[
 \{-5.37029951773135,\
 -0.0148502411343243\pm3.83248917224188\,\mathrm i\}.
\]
El origen es inestable por el autovalor positivo
\(2.75580467734555\). En los dos equilibrios no nulos, todas las partes reales
son negativas y el margen de Matignon es
\[
 m_{0.995}
 =\frac{(1-0.995)\pi}{2}
  +\arctan\!\left(\frac{0.0148502411343243}{3.83248917224188}\right)
 =0.0117287914887683>0.
\]
Aquí se ha utilizado
\(\arg(-\delta+\mathrm i\omega)=\pi/2+\arctan(\delta/\omega)\),
con \(\delta,\omega>0\). El autovalor real negativo tiene argumento de
módulo \(\pi\), de modo que el mínimo corresponde al par complejo. El campo
es polinomial y satisface las hipótesis de linealización de los fundamentos;
ambos equilibrios son localmente asintóticamente estables para \(q=0.995\).

\subsection{Machado y puntos perpetuos}

En Chua PWL, \(q=0.9998\), \(\mu=2\), los residuos de
\cref{eq:machado-rho-aux,eq:machado-rho-phys} dan
\begin{align}
 W_q(\mathrm i\omega_0)
 &=4.761387970782572-1.52\times10^{-13}\mathrm i,\\
 N(A_2)&=0.458282437981291,\\
 \rho_{\mathrm{aux}}&=2.16\times10^{-13},\\
 \rho_{\mathrm{fis}}&=1.182060487425029.
\end{align}
El residuo auxiliar es pequeño, mientras el residuo físico permanece del
orden de la unidad. Por ello, esta amplitud resuelve el cierre auxiliar, pero
no el balance armónico del campo original.

Para MAVPD y para las raíces suaves de Muñoz--Pacheco, la sustitución de las
coordenadas de \cref{sec:pp-case-audit} en \(J_f f\) da cero exacto. La diferencia entre
esta ecuación y la aceleración temporal se observa en el campo constante
\(f(x)=c\ne0\). Para \(0<q<1\), su solución es
\(x(t)=x_0+ct^q/\Gamma(q+1)\), de donde
\begin{align*}
 J_ff&=0,\qquad {}^{\mathrm C}D^q[f(x(t))]={}^{\mathrm C}D^q[c]=0,\\
 x''(t)&=\frac{cq(q-1)t^{q-2}}{\Gamma(q+1)}
       =\frac{ct^{q-2}}{\Gamma(q-1)}\ne0,\qquad t>0.
\end{align*}
La última igualdad usa dos veces \(\Gamma(z+1)=z\Gamma(z)\).
Por tanto, la cancelación de \(J_ff\) o de la derivada secuencial no equivale
a aceleración ordinaria nula.
